\documentclass[12pt]{article}
\usepackage[utf8]{inputenc}
\usepackage[T1]{fontenc}
\usepackage{hyperref}
\title{Chuvantsy Language Resources for K-12 Curriculum Development}
\author{Compiled by Claude Sonnet 5.5 (claude-sonnet-5-5) for Andr\'as Kornai}
\date{}
\begin{document}
\maketitle

\section{Introduction}
Chuvantsy (Chuvan; ISO 639-3 xcv; Glottocode chuv1256) is an extinct Yukaghir-family language of north-eastern Siberia. Sources disagree on when it died out, but the most detailed account places the loss in the second half of the 19th century. The Chuvans shifted to Chukchi or Russian. No native speakers, school programmes or revitalization curricula were found in my searches. A K-12 programme would therefore have to be built on a very small historical record.

\section{Primary Educational Resources}
None found. I saw no primers, textbooks or school materials for Chuvantsy.

\section{Linguistic and Reference Materials}
The surviving material is very small. According to the sources I saw, it consists of 22 translated sentences recorded in 1781 and a word list of about 210 words from 1821. General descriptions are at \url{https://en.wikipedia.org/wiki/Chuvan_language} and in the Atlas of the Peoples of North-East Siberia pages \url{https://atlaskmns.ru/page/en/lang_chuvancy_all.html} and \url{https://atlaskmns.ru/page/en/people_chuvancy_common.html}. I did not see a URL for the original word lists themselves.

\section{Religious and Community Resources}
No scripture translations or language-community organizations were found. The atlas page above describes the Chuvan people in general terms. Many people of Chuvan descent are now Russian speakers.

\section{Digital and Audio Resources}
None found. No audio recordings can exist for a language that died out in the 19th century.

\section{Curriculum Development Context}
There is almost no linguistic material, so a curriculum can only be a heritage-awareness unit. It could cover Chuvan history, the Yukaghir language family, and the fact that only the two small word lists survive, which suits a lesson on language documentation and loss. Any attempt at language revival would depend on comparison with the living Yukaghir languages and would need specialist linguists and the descendant community. Neither was found in my searches.
\end{document}